\begin{definition}
Fix a graph (G), a color bound (k), a distinguished edge (e_0), and a
finite number (r) of noninitial iterations.  The indexed data of a proposed
Tashkinov series consist of tree-sequences (T_i) for
\(0\leq i\leq r+1\), colorings \(\varphi_i\), exact finite color sets
\(S_i\), exact finite edge-identity sets \(F_i\), and optional extension
kinds \(\Theta_i\) for \(0\leq i\leq r\).  They also retain the witnesses
\(v_i,u_i,\delta_i,\gamma_i\) when those are present.

The indexing is zero-based but otherwise literal: Lean index (i) denotes
the manuscript subscript (i).  The sole auxiliary entry (T_0), which is
not one of the manuscript's output tuples, is an auxiliary slot intended
to hold the seed used when certifying the initial closure.  This raw-data
structure by itself imposes no condition on the length of (T_0) or on its
relation to (T_1); the one-edge seed condition and the requirement that
(T_1) be its closure belong to the later Algorithm 3.1 step certificate.
For (i\geq1), the entry
\((T_i,\varphi_{i-1},S_{i-1},F_{i-1},\Theta_{i-1})\) is exactly the
manuscript's (i)-th tuple.  Edge fields store edge identities and therefore
retain multiplicity.
\end{definition}
